\section{Related Work}
\label{sec:related-work}

\subsection{Class-Incremental Learning}
\label{subsec:related-cil}

Class-incremental learning has been approached through parameter
regularization, constrained optimization, distillation, and replay.  Elastic
Weight Consolidation and Synaptic Intelligence protect parameters estimated to
be important for prior tasks~\cite{kirkpatrick2017ewc,zenke2017si}.  GEM and
A-GEM constrain updates using episodic-memory gradients~\cite{lopezpaz2017gem,chaudhry2019agem}.
Learning without Forgetting preserves old outputs by distillation without
retaining old examples~\cite{li2016lwf}.  Uniform Herding belongs to the
replay-based setting and focuses on how exemplar refresh interacts with the
active replay budget.

\subsection{Replay and Exemplar Selection}
\label{subsec:related-replay}

Experience replay retains a bounded subset of prior data and interleaves it
with new-task training~\cite{rolnick2019er,chaudhry2019tiny}.  iCaRL provides
an influential class-incremental exemplar template: it uses herding to select
prioritized exemplars, replays them during training, and predicts with a
nearest-mean-of-exemplars rule~\cite{rebuffi2017icarl}.  The greedy herding
rule approximates a class mean with selected samples~\cite{welling2009herding}.
Other methods formulate selection around gradient diversity or expected
interference~\cite{aljundi2019gss,aljundi2019mir}.

Uniform Herding shares iCaRL's fixed-total active exemplar allocation and greedy
selection primitive, but changes when old classes are refreshed.  After each task it
reselects every observed class in the current representation from its bounded
candidate pool.  iCaRL instead herds a class on arrival and later retains a
prefix of that prioritized list as the quota decreases.  The comparison in
this paper is consequently between complete replay protocols; it does not
isolate refresh while holding every other method choice fixed.

\subsection{Distillation, Classifiers, and Readouts}
\label{subsec:related-readout}

Distillation is frequently paired with replay to preserve old responses.
iCaRL combines output distillation with exemplar replay and NME
classification~\cite{rebuffi2017icarl}; Uniform Herding uses cross-entropy
with temperature-scaled softmax KL on old classes.  The two objectives are
related but not identical: iCaRL's original old-class targets use sigmoid
binary cross-entropy, whereas Uniform Herding uses a softmax KL term.

Classifier design is another source of incremental bias.  LUCIR combines a
normalized cosine classifier with a margin-based incremental objective~\cite{hou2019lucir};
BiC and Weight Aligning explicitly correct or align learned classifier
outputs~\cite{wu2019bic,zhao2020wa}.  Prototype-based classifiers provide a
separate readout: iCaRL predicts from exemplar means, while FeCAM adds
covariance structure to class prototypes~\cite{rebuffi2017icarl,goswami2023fecam}.
Our within-method comparisons separate NME prediction, head-logit prediction,
distillation, and head geometry under the stated training protocol.

\subsection{Forgetting and Stability Metrics}
\label{subsec:related-metrics}

The accuracy matrix used in continual learning supports average accuracy and
backward transfer, as formalized with GEM~\cite{lopezpaz2017gem}.  Average
forgetting summarizes the decline of each old task from its best observed
accuracy to its final accuracy~\cite{chaudhry2019tiny}.  We report all three
quantities so that final accuracy can be read together with retention and
backward transfer.
